\documentclass[11pt,a4paper]{article}

% ---------- Fonts & language ----------
\usepackage[T1]{fontenc}
\usepackage[utf8]{inputenc}
\usepackage[brazil]{babel}
\shorthandoff{"}
\usepackage{tgtermes}
\usepackage{tgadventor}
\usepackage{tgcursor}
\usepackage{listings}

% ---------- Layout ----------
\usepackage[a4paper,top=2.6cm,bottom=2.8cm,left=2.4cm,right=2.4cm]{geometry}
\usepackage{tikz}
\usetikzlibrary{calc,positioning}
\usepackage{xcolor}
\usepackage{booktabs}
\usepackage{tabularx}
\usepackage{enumitem}
\usepackage{amsmath}
\usepackage{amssymb}
\usepackage{tcolorbox}
\tcbuselibrary{skins,breakable}
\usepackage{titlesec}
\usepackage{fancyhdr}
\usepackage[hidelinks,colorlinks=true,linkcolor=accentblue,urlcolor=accentcyan,citecolor=accentblue]{hyperref}
\usepackage{parskip}
\usepackage{microtype}

% ---------- Palette (AI theme) ----------
\definecolor{bgdark}{HTML}{0A0E1A}
\definecolor{bgdark2}{HTML}{131A33}
\definecolor{accentcyan}{HTML}{22D3EE}
\definecolor{accentblue}{HTML}{3B82F6}
\definecolor{accentpurple}{HTML}{8B5CF6}
\definecolor{textlight}{HTML}{E5EDFF}
\definecolor{textmuted}{HTML}{93A4C3}
\definecolor{inkdark}{HTML}{111827}
\definecolor{panelbg}{HTML}{F3F6FD}
\definecolor{panelframe}{HTML}{C7D2FE}

% ---------- Section styling ----------
\titleformat{\section}
  {\sffamily\Large\bfseries\color{inkdark}}
  {\textcolor{accentblue}{\thesection}}{0.6em}{}
\titlespacing*{\section}{0pt}{1.6em}{0.9em}

\titleformat{\subsection}
  {\sffamily\large\bfseries\color{accentblue}}
  {}{0em}{}
\titlespacing*{\subsection}{0pt}{1.3em}{0.5em}

\pagestyle{fancy}
\fancyhf{}
\renewcommand{\headrulewidth}{0.4pt}
\fancyhead[L]{\small\sffamily\color{textmuted} Glossário Técnico --- Projeto LLM}
\fancyhead[R]{\small\sffamily\color{textmuted} \leftmark}
\fancyfoot[C]{\small\sffamily\color{textmuted} \thepage}
\renewcommand{\sectionmark}[1]{\markboth{#1}{}}

\setlist[itemize]{leftmargin=1.3em,itemsep=0.25em}
\setlist[enumerate]{leftmargin=1.5em,itemsep=0.3em}

% ---------- Glossary entry box ----------
\newtcolorbox{termo}[2]{
  breakable,
  enhanced,
  colback=panelbg,
  colframe=panelframe,
  boxrule=0.5pt,
  left=10pt,right=10pt,top=7pt,bottom=7pt,
  arc=3pt,
  title={#1 \hfill \normalfont\small\textcolor{accentblue}{#2}},
  fonttitle=\sffamily\bfseries\color{inkdark},
  coltitle=inkdark,
  titlerule=0.4pt,
  before upper={\parindent0pt},
}

\newcommand{\campo}[2]{\textbf{\textcolor{accentblue}{#1:}} #2\par\smallskip}

\newtcolorbox{notebox}[1][Observação]{
  breakable,
  enhanced,
  colback=white,
  colframe=accentpurple,
  boxrule=0.6pt,
  left=10pt,right=10pt,top=7pt,bottom=7pt,
  arc=2pt,
  borderline west={2.4pt}{0pt}{accentpurple},
  boxrule=0pt,
  title={#1},
  fonttitle=\sffamily\bfseries\small\color{accentpurple},
  coltitle=accentpurple,
  attach title to upper={\par\nobreak\smallskip},
}

\lstdefinestyle{pseudo}{
  basicstyle=\ttfamily\footnotesize\color{textlight},
  backgroundcolor=\color{inkdark},
  frame=single,
  framerule=0pt,
  rulecolor=\color{inkdark},
  framesep=0pt,
  xleftmargin=10pt,
  xrightmargin=10pt,
  framexleftmargin=10pt,
  framexrightmargin=10pt,
  framextopmargin=8pt,
  framexbottommargin=8pt,
  aboveskip=10pt,
  belowskip=10pt,
  showstringspaces=false,
  breaklines=false,
  columns=fullflexible,
  keepspaces=true,
  literate=%
    {á}{{\'a}}1 {à}{{\`a}}1 {ã}{{\~a}}1 {â}{{\^a}}1
    {é}{{\'e}}1 {ê}{{\^e}}1 {í}{{\'i}}1
    {ó}{{\'o}}1 {õ}{{\~o}}1 {ô}{{\^o}}1
    {ú}{{\'u}}1 {ü}{{\"u}}1 {ç}{{\c c}}1
    {Á}{{\'A}}1 {É}{{\'E}}1 {Í}{{\'I}}1 {Ó}{{\'O}}1 {Ú}{{\'U}}1
    {Ã}{{\~A}}1 {Õ}{{\~O}}1 {Â}{{\^A}}1 {Ê}{{\^E}}1 {Ô}{{\^O}}1 {Ç}{{\c C}}1,
}
\lstnewenvironment{pseudocode}{\lstset{style=pseudo}}{}

\begin{document}

% =====================================================
% CAPA
% =====================================================
\begin{titlepage}
\newgeometry{margin=0cm}
\thispagestyle{empty}
\begin{tikzpicture}[remember picture,overlay]
  \shade[top color=bgdark2, bottom color=bgdark]
    (current page.north west) rectangle (current page.south east);

  \begin{scope}
    \pgfmathsetseed{22}
    \foreach \i in {1,...,42}{
      \pgfmathsetmacro{\x}{rnd*21-0.5}
      \pgfmathsetmacro{\y}{rnd*29.7-0.5}
      \coordinate (n\i) at ($(current page.south west)+(\x,\y)$);
    }
    \foreach \k in {1,...,60}{
      \pgfmathtruncatemacro{\ii}{1+mod(int(rnd*10007),42)}
      \pgfmathtruncatemacro{\jj}{1+mod(int(rnd*10007),42)}
      \draw[opacity=0.14,accentcyan,line width=0.3pt] (n\ii) -- (n\jj);
    }
  \end{scope}
  \shade[top color=bgdark2,bottom color=bgdark,opacity=0.55]
    (current page.north west) rectangle ([yshift=8cm]current page.south east);
  \foreach \i in {1,...,42}{
    \fill[accentcyan,opacity=0.55] (n\i) circle (0.6pt);
  }

  \begin{scope}[shift={($(current page.north)+(0,-3.4)$)}]
    \draw[accentcyan,line width=0.9pt] (0,0) -- (0.95,0.55);
    \draw[accentcyan,line width=0.9pt] (0,0) -- (0.95,-0.55);
    \draw[accentpurple,line width=0.9pt] (0,0) -- (-0.95,0.55);
    \draw[accentpurple,line width=0.9pt] (0,0) -- (-0.95,-0.55);
    \draw[accentblue,line width=0.9pt] (0,0) -- (0,1.05);
    \fill[white] (0,0) circle (3.6pt);
    \fill[accentcyan] (0.95,0.55) circle (2.6pt);
    \fill[accentcyan] (0.95,-0.55) circle (2.6pt);
    \fill[accentpurple] (-0.95,0.55) circle (2.6pt);
    \fill[accentpurple] (-0.95,-0.55) circle (2.6pt);
    \fill[accentblue] (0,1.05) circle (2.6pt);
  \end{scope}

  \node[anchor=north] at ($(current page.north)+(0,-5.4)$)
    {\sffamily\color{accentcyan}\fontsize{11}{11}\selectfont
     \MakeUppercase{Glossário Técnico \textbullet\ Sprint 2 \textbullet\ Capítulo 2}};

  \node[anchor=north,align=center,text width=17cm] at ($(current page.north)+(0,-7.1)$)
    {\sffamily\bfseries\color{white}\fontsize{36}{42}\selectfont
     Working with Text Data};

  \node[anchor=north,align=center,text width=15cm] at ($(current page.north)+(0,-10.2)$)
    {\sffamily\color{textlight}\fontsize{15}{20}\selectfont
     De texto bruto a Token IDs, embeddings e posições:\\ o pré-processamento de dados de um LLM};

  \draw[accentcyan,line width=1pt]
    ($(current page.north)+(-2.1,-11.7)$) -- ($(current page.north)+(2.1,-11.7)$);

  \node[anchor=north,align=center,text width=17cm] at ($(current page.north)+(0,-12.8)$)
    {\sffamily\color{textmuted}\fontsize{10}{14}\selectfont
     TOKENIZAÇÃO \ \textbullet\ \ VOCABULÁRIO \ \textbullet\ \ TOKEN IDS \ \textbullet\ \ EMBEDDINGS \ \textbullet\ \ DATALOADER};

  \draw[textmuted,opacity=0.35,line width=0.5pt]
    ($(current page.south west)+(2.2,4.6)$) -- ($(current page.south east)+(-2.2,4.6)$);

  \node[anchor=west] at ($(current page.south west)+(2.2,3.35)$)
    {\sffamily\color{white}\fontsize{12}{14}\selectfont\bfseries Matheus Dapper};
  \node[anchor=west] at ($(current page.south west)+(2.2,2.7)$)
    {\sffamily\color{white}\fontsize{12}{14}\selectfont\bfseries Vitória Aparecida Vendausen};

  \node[anchor=east] at ($(current page.south east)+(-2.2,3.7)$)
    {\sffamily\color{white}\fontsize{12}{14}\selectfont\bfseries 2026};
  \node[anchor=east] at ($(current page.south east)+(-2.2,3.05)$)
    {\sffamily\color{textmuted}\fontsize{9.5}{12}\selectfont Sprint 2 --- 27 de agosto};

\end{tikzpicture}
\restoregeometry
\end{titlepage}

\pagestyle{empty}
\tableofcontents
\clearpage
\pagestyle{fancy}
\setcounter{page}{1}

% =====================================================
\section*{Introdução}
\addcontentsline{toc}{section}{Introdução}
\markboth{Introdução}{}

Este glossário registra os conceitos introduzidos no Capítulo 2 --- \emph{Working with
Text Data} --- do livro \emph{Build a Large Language Model (From Scratch)}
(Sebastian Raschka, Manning, 2025), referência principal do Projeto LLM. O objetivo não
é traduzir os termos da língua inglesa, mas registrar de forma sintética seu
significado e sua função no pipeline de pré-processamento de dados de um modelo de
linguagem. O glossário é cumulativo em relação ao produzido na Sprint 1 (ver
\texttt{LLM\_Glossario\_IA.tex}) e deverá ser reutilizado nas próximas Sprints.

% =====================================================
\section{Tokenização e tokens}

\begin{termo}{Tokenization}{Tokenização}
\campo{Definição}{Processo de dividir uma sequência contínua de texto em unidades
menores e discretas --- os tokens --- que podem ser palavras, partes de palavras,
sinais de pontuação ou até caracteres/bytes individuais, dependendo do esquema
utilizado.}
\campo{Função no modelo}{É a primeira etapa do pipeline: transforma um objeto que a
rede não sabe processar (uma \emph{string}) em uma sequência de unidades discretas que
podem, na etapa seguinte, ser mapeadas para números.}
\campo{Relação com outros conceitos}{Precede a construção do \emph{Vocabulário} e a
conversão para \emph{Token ID}. No GPT, a tokenização é feita por
\emph{Byte Pair Encoding (BPE)}, não pela versão simples baseada em regex.}
\campo{Exemplo}{\texttt{"Hello, do you like tea?"} $\rightarrow$ \texttt{["Hello", ",",
"do", "you", "like", "tea", "?"]} (tokenização por regex, separando palavras e
pontuação).}
\end{termo}

\begin{termo}{Token}{Token}
\campo{Definição}{Unidade mínima produzida pela tokenização. Pode corresponder a uma
palavra inteira, a um pedaço de palavra (subword) ou a um sinal de pontuação --- nunca a
um conceito semântico completo por si só.}
\campo{Função no modelo}{É a unidade de trabalho de todo o pipeline: o texto deixa de
ser manipulado como caracteres soltos e passa a ser manipulado como uma sequência de
tokens, cada um posteriormente associado a um \emph{Token ID} e a um vetor de
\emph{Embedding (representação vetorial)}.}
\campo{Relação com outros conceitos}{Um token não carrega significado por si --- seu
``significado'' computacional só surge depois, quando seu \emph{Token ID} é
convertido em vetor pela camada de embedding.}
\campo{Exemplo}{Na frase ``the sunlit terraces'', o BPE do GPT-2 produz os tokens
\texttt{["the", " sun", "lit", " terr", "aces"]} --- note que ``sunlit'' e ``terraces''
são quebrados em subpalavras, pois não aparecem inteiras no vocabulário do
tokenizador.}
\end{termo}

\begin{termo}{Corpus}{Corpus (de texto)}
\campo{Definição}{Conjunto de texto bruto utilizado como fonte de dados --- para a
construção do vocabulário, para a extração de sequências de treinamento, ou para o
pré-treinamento do modelo.}
\campo{Função no modelo}{É a matéria-prima de toda a Sprint: sem um corpus não há
tokens, não há vocabulário e não há sequências de treinamento.}
\campo{Relação com outros conceitos}{Alimenta diretamente a \emph{Tokenization}; seu
tamanho (em tokens) determina quantas amostras de \emph{Sequência de treinamento
(par entrada/alvo)} podem ser extraídas para um dado tamanho de contexto.}
\campo{Exemplo}{Neste projeto, utiliza-se o texto \emph{``The Verdict''}
(Edith Wharton, domínio público), o mesmo corpus de exemplo do capítulo, com
aproximadamente 20\,479 caracteres e 5\,145 tokens BPE.}
\end{termo}

% =====================================================
\section{Vocabulário e Token IDs}

\begin{termo}{Vocabulary}{Vocabulário}
\campo{Definição}{Conjunto de todos os tokens distintos que um tokenizador é capaz de
reconhecer, organizado como um dicionário que associa cada token a um número inteiro
único (seu \emph{Token ID}).}
\campo{Função no modelo}{Define o espaço discreto sobre o qual o modelo opera: todo
token de entrada precisa existir no vocabulário (ou ser tratado por um token especial)
para ser convertido em ID, e todo token gerado pelo modelo na saída é, na verdade, um
índice desse mesmo vocabulário.}
\campo{Relação com outros conceitos}{Construído a partir dos tokens únicos de um
\emph{Corpus}; define o tamanho da tabela de \emph{Token embedding}
(\texttt{vocab\_size} linhas) e do vetor de \emph{logits} produzido na camada final do
modelo.}
\campo{Exemplo}{No corpus \emph{The Verdict}, o vocabulário construído a partir dos
tokens únicos possui 1\,130 entradas; somando os dois tokens especiais
(\texttt{<|unk|>} e \texttt{<|endoftext|>}), chega a 1\,132. Já o vocabulário BPE do
GPT-2, fixo e independente do corpus, possui 50\,257 entradas.}
\end{termo}

\begin{termo}{Token ID}{Token ID (identificador do token)}
\campo{Definição}{Número inteiro que representa um token dentro do vocabulário --- o
índice da linha correspondente a esse token na tabela de vocabulário.}
\campo{Função no modelo}{É a única representação numérica bruta do texto que a rede
manipula antes da camada de embedding. Serve como índice de consulta (\emph{lookup}) na
tabela de \emph{Token embedding}.}
\campo{Relação com outros conceitos}{Obtido a partir de um \emph{Token} via o
\emph{Vocabulary}; convertido de volta em token pela operação inversa
(\emph{decoding}); não deve ser confundido com uma representação semântica --- dois IDs
numericamente próximos (ex.: 5 e 6) não implicam tokens semanticamente relacionados.}
\campo{Exemplo}{\texttt{"Hello"} $\rightarrow$ Token ID \texttt{15496} no vocabulário
BPE do GPT-2 (\texttt{tiktoken}, \texttt{encoding="gpt2"}).}
\end{termo}

\begin{termo}{Special tokens}{Tokens especiais}
\campo{Definição}{Tokens artificiais, sem correspondência direta no texto original,
adicionados ao vocabulário para tratar casos que a tokenização normal não resolve. Os
dois discutidos no capítulo são \texttt{<|unk|>} (\emph{unknown}, para palavras fora do
vocabulário) e \texttt{<|endoftext|>} (para marcar a fronteira entre documentos
concatenados).}
\campo{Função no modelo}{\texttt{<|unk|>} evita que o \emph{Tokenization} falhe
diante de uma palavra nunca vista; \texttt{<|endoftext|>} evita que o modelo aprenda
associações espúrias entre o fim de um texto e o início de outro quando vários
documentos são concatenados em um único fluxo de treinamento.}
\campo{Relação com outros conceitos}{São entradas adicionais no \emph{Vocabulary},
recebendo seus próprios \emph{Token ID}; o tokenizador \emph{Byte Pair Encoding
(BPE)} do GPT, por lidar com qualquer string via subpalavras, normalmente dispensa o
\texttt{<|unk|>} --- mas mantém o \texttt{<|endoftext|>}.}
\campo{Exemplo}{\texttt{"Hello, do you like tea?"} contra um vocabulário que não conhece
``Hello'' produz \texttt{["{}<|unk|>", ",", "do", "you", "like", "tea", "?"]}.}
\end{termo}

\begin{termo}{Byte Pair Encoding (BPE)}{Codificação por Pares de Bytes}
\campo{Definição}{Algoritmo de tokenização em nível de subpalavra que constrói seu
vocabulário mesclando iterativamente os pares de caracteres (ou bytes) mais frequentes
em um corpus de treinamento, até atingir um tamanho de vocabulário alvo.}
\campo{Função no modelo}{É o tokenizador efetivamente usado pelo GPT. Por operar em
subpalavras e, no limite, em bytes individuais, consegue representar qualquer string de
entrada --- inclusive palavras inventadas ou fora do vocabulário --- sem nunca precisar
de um token \texttt{<|unk|>}.}
\campo{Relação com outros conceitos}{Substitui, na prática, o
\emph{Tokenization} + \emph{Vocabulary} caseiros descritos nas seções iniciais
do capítulo; implementado neste projeto via a biblioteca \texttt{tiktoken}
(\texttt{encoding="gpt2"}), com vocabulário fixo de 50\,257 tokens.}
\campo{Exemplo}{A palavra inventada \texttt{"Akwirw ier"} é quebrada pelo BPE em
\texttt{["Ak", "w", "ir", "w", " ", "ier"]} --- seis subpalavras que, recombinadas,
reconstroem exatamente a string original.}
\end{termo}

% =====================================================
\clearpage
\section{Sequências de treinamento}

\begin{termo}{Sliding window}{Janela deslizante}
\campo{Definição}{Técnica de amostragem que percorre a sequência de Token IDs do corpus
com uma janela de tamanho fixo (\texttt{max\_length}), avançando um número fixo de
posições (\texttt{stride}) a cada passo, para extrair múltiplas subsequências de
treinamento a partir de um único corpus contínuo.}
\campo{Função no modelo}{É o mecanismo que transforma um corpus (uma longa sequência de
tokens) em um conjunto de exemplos de treinamento de tamanho fixo, prontos para compor
lotes (\emph{batches}).}
\campo{Relação com outros conceitos}{Produz os pares de \emph{Sequência de
treinamento (par entrada/alvo)}; o \texttt{stride} controla a sobreposição entre
janelas consecutivas e, junto com o \emph{Context length (tamanho do contexto)},
determina a quantidade total de amostras geradas.}
\campo{Exemplo}{Com \texttt{max\_length=4} e \texttt{stride=4} (sem sobreposição), o
corpus \emph{The Verdict} (5\,145 tokens BPE) produz 1\,286 amostras; com
\texttt{stride=1} (máxima sobreposição), produziria cerca de 5\,141 amostras.}
\end{termo}

\begin{termo}{Input--target pair}{Par entrada/alvo (sequência de treinamento)}
\campo{Definição}{Par de sequências de mesmo tamanho em que o alvo (\emph{target}) é
exatamente a sequência de entrada (\emph{input}) deslocada em uma posição para a
direita.}
\campo{Função no modelo}{Implementa, de forma totalmente automática, a tarefa de
auto-supervisão do LLM: ``prever o próximo token''. Não é necessária nenhuma rotulação
manual --- o próprio texto fornece o rótulo de cada posição.}
\campo{Relação com outros conceitos}{Gerado pela \emph{Sliding window} sobre os
\emph{Token ID} do corpus; consumido pelo \emph{DataLoader}, que os agrupa em
lotes (\emph{Batch (lote)}) para alimentar o modelo durante o treinamento.}
\campo{Exemplo}{Para a sequência de IDs \texttt{[40, 367, 2885, 1464]}: entrada =
\texttt{[40, 367, 2885]}, alvo = \texttt{[367, 2885, 1464]}. Em texto:
entrada = \texttt{"I HAD always"}, alvo = \texttt{" HAD always thought"}.}
\end{termo}

\begin{termo}{Context length}{Tamanho do contexto (janela de contexto)}
\campo{Definição}{Número máximo de tokens que uma amostra de entrada contém
simultaneamente --- e, por extensão, o número máximo de tokens anteriores que o modelo
consegue considerar ao prever o próximo.}
\campo{Função no modelo}{Determina a dimensão da sequência de entrada e da tabela de
\emph{Positional embedding}; limita quanto contexto passado o modelo pode usar em
cada predição.}
\campo{Relação com outros conceitos}{Parâmetro central da \emph{Sliding window};
quanto maior o contexto, menos amostras um corpus de tamanho fixo produz (quando o
\texttt{stride} acompanha o contexto), mas mais informação cada amostra carrega.}
\campo{Exemplo}{Contexto de 4 tokens gera 1\,286 amostras no corpus de teste; contexto
de 128 tokens gera apenas 40 amostras a partir do mesmo corpus (ver Seção de
Experimentos).}
\end{termo}

\begin{termo}{Dataset}{Conjunto de dados (\texttt{torch.utils.data.Dataset})}
\campo{Definição}{Estrutura que encapsula a lógica de acesso a um conjunto de amostras
de treinamento --- no caso deste projeto, a classe \texttt{GPTDatasetV1}, que aplica a
\emph{Sliding window} sobre o corpus tokenizado e expõe cada par (entrada, alvo)
por índice.}
\campo{Função no modelo}{Padroniza a forma como os dados são armazenados e acessados,
permitindo que o \emph{DataLoader} os agrupe em lotes de forma genérica, sem
conhecer os detalhes do corpus.}
\campo{Relação com outros conceitos}{Implementa a interface exigida pelo
\emph{DataLoader} do PyTorch (\texttt{\_\_len\_\_} e \texttt{\_\_getitem\_\_});
armazena internamente os \emph{Input--target pair} pré-computados.}
\campo{Exemplo}{\texttt{GPTDatasetV1(text, tokenizer, max\_length=8, stride=8)} produz
um objeto com \texttt{len(dataset) == 643} amostras para o corpus de teste.}
\end{termo}

\begin{termo}{DataLoader}{DataLoader (carregador de dados)}
\campo{Definição}{Componente do PyTorch (\texttt{torch.utils.data.DataLoader}) que
envolve um \emph{Dataset} e organiza suas amostras em lotes (\emph{batches}),
podendo embaralhar (\emph{shuffle}) a ordem das amostras a cada época e paralelizar o
carregamento.}
\campo{Função no modelo}{É a interface final entre os dados preparados e o laço de
treinamento: a cada iteração, entrega um lote de tensores \texttt{(entrada, alvo)}
prontos para ser processados pela rede.}
\campo{Relação com outros conceitos}{Consome um \emph{Dataset} (aqui,
\texttt{GPTDatasetV1}); seu parâmetro \texttt{batch\_size} define o tamanho do
\emph{Batch (lote)}, e \texttt{drop\_last} controla se o último lote incompleto de
uma época é descartado.}
\campo{Exemplo}{\texttt{create\_dataloader\_v1(text, batch\_size=8, max\_length=4,
stride=4)} produz lotes com \texttt{inputs.shape == (8, 4)}.}
\end{termo}

\begin{termo}{Batch}{Lote}
\campo{Definição}{Conjunto de \texttt{batch\_size} amostras de entrada (e seus
respectivos alvos) agrupadas em um único tensor, processadas simultaneamente por uma
passagem da rede.}
\campo{Função no modelo}{Permite explorar o paralelismo de hardware (GPU/CPU) e
estabiliza a estimativa do gradiente durante o treinamento, ao calcular a perda sobre
várias amostras de uma vez em vez de uma única amostra isolada.}
\campo{Relação com outros conceitos}{Produzido pelo \emph{DataLoader}; sua
dimensão (\texttt{batch\_size}) é o primeiro eixo dos tensores de
\emph{Token embedding} e \emph{Input embedding}, de forma
\texttt{(batch\_size, context\_length, output\_dim)}.}
\campo{Exemplo}{Com \texttt{batch\_size=8} e \texttt{max\_length=4}, cada lote de
entrada tem forma \texttt{(8, 4)}; após a camada de embedding com
\texttt{output\_dim=256}, passa a ter forma \texttt{(8, 4, 256)}.}
\end{termo}

% =====================================================
\clearpage
\section{Embeddings e posição}

\begin{termo}{Embedding}{Embedding (representação vetorial densa)}
\campo{Definição}{Vetor de números reais, de dimensão fixa (\texttt{output\_dim}), que
representa um objeto discreto (aqui, um token) em um espaço vetorial contínuo. Ao
contrário do Token ID, os valores do vetor são aprendidos durante o treinamento.}
\campo{Função no modelo}{Permite que a rede opere sobre representações contínuas e
diferenciáveis, sobre as quais operações como produto interno e multiplicação de
matrizes fazem sentido geométrico --- o que não é possível com um índice inteiro
isolado.}
\campo{Relação com outros conceitos}{Produzido a partir de um \emph{Token ID} por
uma \emph{Embedding layer}; a soma do \emph{Token embedding} com o
\emph{Positional embedding} resulta no \emph{Input embedding}.}
\campo{Exemplo}{Com \texttt{output\_dim=256}, cada Token ID é mapeado para um vetor de
256 números reais --- por exemplo, o token de ID 40 (``I'') passa a ser representado
por um vetor \texttt{[0.34, -1.19, \ldots]} (256 valores) em vez do escalar
\texttt{40}.}
\end{termo}

\begin{termo}{Embedding layer}{Camada de embedding (\texttt{nn.Embedding})}
\campo{Definição}{Camada de rede neural que funciona como uma tabela de consulta
(\emph{lookup table}) de dimensão \texttt{(vocab\_size, output\_dim)}: dado um Token
ID, retorna a linha correspondente da tabela. Os valores da tabela são parâmetros
treináveis, inicializados aleatoriamente.}
\campo{Função no modelo}{É a ponte entre o mundo discreto dos \emph{Token ID} e o
mundo contínuo em que a rede neural opera. Equivale, matematicamente, a multiplicar um
vetor \emph{one-hot} pela matriz de pesos da camada --- mas implementada de forma muito
mais eficiente, como indexação direta.}
\campo{Relação com outros conceitos}{Instanciada com \texttt{vocab\_size} igual ao
tamanho do \emph{Vocabulary} do tokenizador utilizado; sua saída, somada à saída de
uma segunda \texttt{nn.Embedding} indexada por posição, forma o
\emph{Input embedding}.}
\campo{Exemplo}{\texttt{nn.Embedding(50257, 256)} cria uma tabela com
$50\,257 \times 256 = 12\,865\,792$ parâmetros treináveis --- um por célula da tabela.}
\end{termo}

\begin{termo}{Positional embedding}{Embedding posicional}
\campo{Definição}{Vetor, de mesma dimensão do \emph{Token embedding}, que codifica
exclusivamente a posição de um token dentro da sequência (0, 1, 2, \ldots), e não seu
conteúdo. Neste projeto, segue a abordagem do GPT: posicionamento absoluto e aprendido,
via uma segunda \texttt{nn.Embedding} indexada pela posição.}
\campo{Função no modelo}{Resolve uma limitação fundamental da camada de embedding
pura: como ela é uma tabela de consulta por Token ID, o mesmo token gera sempre o mesmo
vetor, não importa em que posição da sequência apareça. Sem informação posicional, o
modelo enxergaria a sequência como um conjunto (\emph{bag of tokens}), não como uma
sequência ordenada.}
\campo{Relação com outros conceitos}{Somado ao \emph{Token embedding} para formar o
\emph{Input embedding}; sua tabela tem tamanho
\texttt{(context\_length, output\_dim)}, diferente da tabela de token embeddings, que
tem tamanho \texttt{(vocab\_size, output\_dim)}.}
\campo{Exemplo}{O token de ID 40 na posição 0 e o mesmo token de ID 40 na posição 2
recebem, após a soma posicional, vetores diferentes --- verificado experimentalmente
em \texttt{src/embeddings.py} (\texttt{vetor\_pos0 != vetor\_pos2}).}
\end{termo}

\begin{termo}{Input embedding}{Embedding de entrada (token + posição)}
\campo{Definição}{Soma elemento a elemento entre o \emph{Token embedding} de cada
posição e o \emph{Positional embedding} correspondente àquela posição. É o vetor
final que representa cada token da sequência ao entrar nos blocos do Transformer.}
\campo{Função no modelo}{É literalmente a entrada da arquitetura GPT --- o ponto de
saída desta Sprint e o ponto de entrada da Sprint seguinte (mecanismo de atenção).
Carrega simultaneamente informação de conteúdo (o que o token é) e de posição (onde ele
está na sequência).}
\campo{Relação com outros conceitos}{Resultado de \emph{Token embedding} +
\emph{Positional embedding}; tem forma
\texttt{(batch\_size, context\_length, output\_dim)}, a mesma forma esperada pela
primeira camada de \emph{self-attention} (Capítulo 3).}
\campo{Exemplo}{Para um lote com \texttt{batch\_size=8}, \texttt{context\_length=4} e
\texttt{output\_dim=256}, o \emph{input embedding} tem forma \texttt{(8, 4, 256)}.}
\end{termo}

\end{document}
